\section{Proceso completo del workflow}

\subsection{Proceso de seis pasos}

\begin{enumerate}
    \item \textbf{Paso 1: solución inicial}——produce una solución inicial a partir del prompt y de Gemini 2.5 Pro
    \item \textbf{Paso 2: Self-Improvement}——comparte el contexto del paso 1 y continúa haciendo refine (porque el paso 1 puede terminar de manera apresurada por agotarse el Thinking Budget)
    \item \textbf{Paso 3: Verification}——verifica la solución del paso 2 y produce un verdict (válido/inválido) y un bug report
    \item \textbf{Paso 4: Bug Report}——extrae la ubicación y la causa concretas del error
    \item \textbf{Paso 5: Correction}——corrige la solution con base en el bug report
    \item \textbf{Paso 6: ciclo}——los pasos 3 a 5 se repiten en ciclo
\end{enumerate}

\begin{importantbox}{Condiciones de terminación}
\begin{itemize}
    \item \textbf{Aceptación}: una solution pasa la verification \textbf{5 veces} consecutivas $\rightarrow$ se considera resuelto el problema
    \item \textbf{Rechazo}: tras \textbf{10} ciclos sigue sin pasar $\rightarrow$ se abandona la dirección actual y se vuelve a hacer sample
\end{itemize}
\end{importantbox}

\subsection{Patrón Evaluator-Optimizer}

\begin{knowledgebox}{Correspondencia con el patrón clásico de Building Effective Agents}
Este workflow es, en esencia, el patrón \textbf{Evaluator-Optimizer}:
\begin{itemize}
    \item el Generator (generador) produce la solution
    \item el Evaluator (evaluador) la verifica y da feedback
    \item si se acepta, se produce como salida; si se rechaza, se vuelve a generar con base en el feedback
\end{itemize}
\end{knowledgebox}

\subsection{El papel decisivo de la calidad de la solución inicial}

\begin{warningbox}{La clave del éxito está en la solución inicial}
Si la solución inicial que se genera en el paso 2 tiene un overlap considerable con la solución correcta, entonces el ciclo posterior de verificación y corrección tiene una alta probabilidad de resolver el problema. Pero si la dirección de la solución inicial ya está equivocada (se desvía demasiado), es poco probable que las correcciones posteriores puedan salvarla.
\end{warningbox}

\subsection{Resumen de la sección}

Todo el workflow es un ciclo iterativo estándar de generación, verificación y corrección, cuyas condiciones de terminación se basan en el conteo de veces consecutivas que se pasa o no se pasa la verificación.
